\documentclass{article}
\usepackage{amsmath,amssymb,amsthm}
\usepackage{algorithm,algorithmic}
\usepackage{hyperref}
\usepackage{geometry}
\geometry{margin=1in}
\newtheorem{theorem}{Theorem}
\newtheorem{lemma}{Lemma}
\newtheorem{assumption}{Assumption}
\newcommand{\E}{\mathbb{E}}
\newcommand{\R}{\mathbb{R}}

\begin{document}

\section*{Gradient Descent‑Ascent (GDA) for Non‑Convex‑Concave Min‑Max}

\section*{Mathematical Formulation}
Consider the saddle‑point problem  
\[
\min_{y\in\R^{d_y}}\;\max_{x\in\R^{d_x}} \;L(x,y),
\]
where the objective $L:\R^{d_x}\times\R^{d_y}\rightarrow\R$ is \emph{smooth} (i.e. its gradient is $L$‑Lipschitz) but is **not** assumed convex in $x$ nor concave in $y$.

Define the concatenated variable  
\[
z:=\begin{bmatrix}x\\-y\end{bmatrix}\in\R^{d},\qquad d:=d_x+d_y,
\]
and the \emph{game potential}
\[
\Phi(z):=L\bigl(x,-\,\!(-y)\bigr)=L(x,y).
\]
Its gradient w.r.t.\ $z$ is
\[
\nabla\Phi(z)=\begin{bmatrix}
\nabla_x L(x,y)\\[2mm]
-\nabla_y L(x,y)
\end{bmatrix}
=:F(z).
\]
Hence a single GDA step with step size $\eta>0$ is exactly a gradient–descent step on $\Phi$:
\[
z_{t+1}=z_t-\eta F(z_t)
\quad\Longleftrightarrow\quad
\begin{cases}
x_{t+1}=x_t+\eta \,\nabla_x L(x_t,y_t),\\[1mm]
y_{t+1}=y_t-\eta \,\nabla_y L(x_t,y_t).
\end{cases}
\tag{1}
\]

\section*{Pseudocode}
\begin{algorithm}[H]
\caption{Gradient Descent‑Ascent (GDA)}
\begin{algorithmic}[1]
\REQUIRE Initial pair $(x_0,y_0)$, step size $\eta>0$, number of iterations $T$.
\FOR{$t=0,\dots,T-1$}
    \STATE Compute stochastic (or deterministic) gradients  
    \[
    g^x_t=\nabla_x L(x_t,y_t),\qquad
    g^y_t=\nabla_y L(x_t,y_t).
    \]
    \STATE Update the primal (max) variable (ascent)  
    \[
    x_{t+1}=x_t+\eta\, g^x_t .
    \]
    \STATE Update the dual (min) variable (descent)  
    \[
    y_{t+1}=y_t-\eta\, g^y_t .
    \]
\ENDFOR
\RETURN $(x_T,y_T)$.
\end{algorithmic}
\end{algorithm}

\section*{Dry Run Example}
We illustrate the update on the simple **single‑variable** function  
\[
f(w)=(w-3)^2 .
\]
Here $L\equiv -f$ (so that we are maximising $L$) and there is no $y$ variable.  
Gradient: $\nabla f(w)=2(w-3)$.  
We initialise $w_0=0$ and use step size $\eta=0.1$ for three iterations.

\begin{center}
\begin{tabular}{c|c|c|c}
$t$ & $w_t$ & $\nabla f(w_t)=2(w_t-3)$ & $w_{t+1}=w_t-\eta\nabla f(w_t)$ \\ \hline
0 & $0$ & $-6$ & $0-0.1(-6)=0.6$ \\[1mm]
1 & $0.6$ & $2(0.6-3)=-4.8$ & $0.6-0.1(-4.8)=1.08$ \\[1mm]
2 & $1.08$ & $2(1.08-3)=-3.84$ & $1.08-0.1(-3.84)=1.464$ \\[1mm]
3 & $1.464$ & $2(1.464-3)=-3.072$ & $1.464-0.1(-3.072)=1.7712$
\end{tabular}
\end{center}
Objective values $f(w_t)$ decrease from $9$ to $3.22$, illustrating the descent direction.

\newpage
\section*{Assumptions}
\begin{assumption}[Smoothness]\label{as:smooth}
There exists $L>0$ such that for all $z,z'\in\R^{d}$,
\[
\|\nabla\Phi(z)-\nabla\Phi(z')\| \le L\|z-z'\|.
\tag{A1}
\]
Equivalently, $\Phi$ has $L$‑Lipschitz gradient.
\end{assumption}

\begin{assumption}[Lower Boundedness]\label{as:lb}
The game potential $\Phi$ is bounded below:
\[
\Phi^{\inf}:=\inf_{z\in\R^{d}}\Phi(z)>-\infty .
\tag{A2}
\end{assumption}

\begin{assumption}[Step Size]\label{as:step}
The constant step size satisfies
\[
0<\eta\le\frac{1}{L}.
\tag{A3}
\end{assumption}

\section*{Main Convergence Theorem}
\begin{theorem}\label{thm:main}
Assume \ref{as:smooth}–\ref{as:step}. Let $\{(x_t,y_t)\}_{t=0}^{T}$ be the iterates generated by \eqref{1}. Then
\[
\frac{1}{T}\sum_{t=0}^{T-1}\bigl\|\nabla_x L(x_t,y_t)\bigr\|^{2}
      +\bigl\|\nabla_y L(x_t,y_t)\bigr\|^{2}
   \;\le\;
   \frac{2\bigl(\Phi(z_0)-\Phi^{\inf}\bigr)}{\eta\,T}.
\tag{2}
\]
Consequently,
\[
\min_{0\le t<T}\Bigl\|\begin{bmatrix}\nabla_x L(x_t,y_t)\\[1mm]-\nabla_y L(x_t,y_t)\end{bmatrix}\Bigr\|
\;=\;
O\!\left(\frac{1}{\sqrt{T}}\right).
\]
\end{theorem}

\section*{Theoretical Properties}
\begin{itemize}
\item \textbf{Stationarity.} The bound (2) guarantees that there exists at least one iterate whose \emph{game gradient}
\(
F(z_t)=\bigl[\nabla_x L(x_t,y_t);\,-\nabla_y L(x_t,y_t)\bigr]
\)
has norm $O(1/\sqrt{T})$, i.e. an \emph{approximate first‑order stationary point} of the min‑max problem.

\item \textbf{Hyper‑parameter sensitivity.}
The step size must obey $\eta\le 1/L$ (Assumption \ref{as:step}). Choosing a larger $\eta$ can break the descent inequality used in the proof and may cause divergence, while a smaller $\eta$ simply slows the $1/\sqrt{T}$ rate.

\item \textbf{Comparison to baselines.}
For convex–concave $L$, extragradient or optimistic GDA achieve $O(1/T)$ in terms of the primal–dual gap. In the non‑convex‑concave regime the guarantee deteriorates to $O(1/\sqrt{T})$ on the gradient norm, which matches the optimal rate for smooth non‑convex optimization (see \cite{ghadimi2013stochastic}).

\item \textbf{Extension to stochastic gradients.}
If $g_t$ is an unbiased estimator of $F(z_t)$ with bounded variance $\sigma^2$, the same proof yields an additional $\eta\sigma^2/T$ term, preserving the $O(1/\sqrt{T})$ rate after the usual choice $\eta=O(1/\sqrt{T})$.
\end{itemize}

\section*{Discussion}
The theorem shows that, despite the lack of convexity/concavity, simple GDA still converges to a stationary point at the same rate as gradient descent on a smooth non‑convex function. The key insight is the re‑parameterisation $z=(x,-y)$ which turns the min‑max update into ordinary gradient descent on the smooth scalar function $\Phi(z)$. This reduction allows us to reuse standard non‑convex analysis. The result is tight in the sense that, for general smooth non‑convex functions, no algorithm based solely on first‑order information can achieve a better worst‑case dependence on $T$.

\newpage
\section*{Preliminaries (Definitions and Lemmas)}
\begin{definition}[L‑smooth function]
A differentiable function $\Phi:\R^{d}\!\to\!\R$ is $L$‑smooth if for all $z,z'$,
\[
\Phi(z')\le \Phi(z)+\langle\nabla\Phi(z),z'-z\rangle+\frac{L}{2}\|z'-z\|^{2}.
\tag{D1}
\]
\end{definition}

\begin{lemma}[Descent Lemma]\label{lem:descent}
If $\Phi$ is $L$‑smooth, then for any $z$ and step size $\eta\le 1/L$,
\[
\Phi(z-\eta\nabla\Phi(z))
\le \Phi(z)-\frac{\eta}{2}\|\nabla\Phi(z)\|^{2}.
\tag{L1}
\]
\end{lemma}
\begin{proof}
Apply (D1) with $z'=z-\eta\nabla\Phi(z)$:
\[
\Phi(z-\eta\nabla\Phi(z))
\le \Phi(z)-\eta\|\nabla\Phi(z)\|^{2}
      +\frac{L\eta^{2}}{2}\|\nabla\Phi(z)\|^{2}.
\]
Since $\eta\le 1/L$, the coefficient of $\|\nabla\Phi(z)\|^{2}$ is at most $-\frac{\eta}{2}$, yielding (L1). \hfill $\square$
\end{proof}

\section*{Main Proof (Step‑by‑Step Derivation)}
We prove Theorem~\ref{thm:main}. Recall $z_t=[x_t;-y_t]$ and $F(z_t)=\nabla\Phi(z_t)$. The GDA update is
\[
z_{t+1}=z_t-\eta F(z_t).
\tag{P1}
\]

\begin{enumerate}
\item[Step 1.] Apply Lemma~\ref{lem:descent} with $z=z_t$ and step size $\eta$ (valid by Assumption \ref{as:step}):
\[
\Phi(z_{t+1})
\le \Phi(z_t)-\frac{\eta}{2}\|F(z_t)\|^{2}.
\tag{S1}
\]

\item[Step 2.] Rearrange (S1) to isolate the squared norm:
\[
\frac{\eta}{2}\|F(z_t)\|^{2}
\le \Phi(z_t)-\Phi(z_{t+1}).
\tag{S2}
\]

\item[Step 3.] Sum (S2) over $t=0,\dots,T-1$:
\[
\frac{\eta}{2}\sum_{t=0}^{T-1}\|F(z_t)\|^{2}
\le \Phi(z_0)-\Phi(z_T).
\tag{S3}
\]

\item[Step 4.] Use the lower bound (A2): $\Phi(z_T)\ge\Phi^{\inf}$, thus
\[
\frac{\eta}{2}\sum_{t=0}^{T-1}\|F(z_t)\|^{2}
\le \Phi(z_0)-\Phi^{\inf}.
\tag{S4}
\]

\item[Step 5.] Divide both sides of (S4) by $\eta T$:
\[
\frac{1}{T}\sum_{t=0}^{T-1}\|F(z_t)\|^{2}
\le \frac{2\bigl(\Phi(z_0)-\Phi^{\inf}\bigr)}{\eta\,T}.
\tag{S5}
\]

\item[Step 6.] Unpack $F(z_t)$:
\[
\|F(z_t)\|^{2}
= \bigl\|\nabla_x L(x_t,y_t)\bigr\|^{2}
  +\bigl\|\nabla_y L(x_t,y_t)\bigr\|^{2}.
\tag{S6}
\]

\item[Step 7.] Substitute (S6) into (S5) to obtain exactly the bound (2) of Theorem~\ref{thm:main}.
\hfill $\square$
\end{enumerate}

\section*{Rate Derivation (With Constants)}
From (S5) we have
\[
\min_{0\le t<T}\|F(z_t)\|
\;\le\;
\sqrt{\frac{2\bigl(\Phi(z_0)-\Phi^{\inf}\bigr)}{\eta\,T}}.
\]
Choosing the largest admissible step size $\eta=1/L$ (Assumption \ref{as:step}) gives the explicit constant
\[
\min_{0\le t<T}\|F(z_t)\|
\;\le\;
\sqrt{\frac{2L\bigl(\Phi(z_0)-\Phi^{\inf}\bigr)}{T}}
\;=\;O\!\left(\frac{1}{\sqrt{T}}\right).
\]

\section*{Special Cases}
\begin{itemize}
\item \textbf{Convex–concave $L$.} If $L$ is convex in $x$ and concave in $y$, the game gradient $F$ is \emph{monotone}. Under the same step size, the bound (2) still holds, but stronger results (primal–dual gap $O(1/T)$) become available via extragradient methods.

\item \textbf{Deterministic vs. stochastic gradients.} The proof above uses the exact gradient $F(z_t)$. When only an unbiased estimator $\tilde F_t$ with $\E[\tilde F_t|z_t]=F(z_t)$ and $\E\|\tilde F_t-F(z_t)\|^{2}\le\sigma^{2}$ is available, the same steps give
\[
\frac{1}{T}\sum_{t=0}^{T-1}\E\|F(z_t)\|^{2}
\le\frac{2(\Phi(z_0)-\Phi^{\inf})}{\eta T}
     +\eta\sigma^{2},
\]
and the optimal choice $\eta=O(1/\sqrt{T})$ yields the $O(1/\sqrt{T})$ rate.

\item \textbf{Strongly‑smooth potential.} If $\Phi$ is $L$‑smooth and also $\mu$‑strongly convex, the classic GD analysis gives a linear rate $\bigl(1-\eta\mu\bigr)^{T}$, which is a strictly stronger guarantee than (2). This situation reduces to the well‑studied convex–concave case.
\end{itemize}

\newpage
\begin{thebibliography}{9}
\bibitem{ghadimi2013stochastic}
S.~Ghadimi and G.~Lan, ``Stochastic First‑and Zero‑Order Methods for Nonconvex Stochastic Programming,'' \emph{SIAM Journal on Optimization}, vol.~23, no.~4, pp.~2341–2368, 2013.
\end{thebibliography}

\end{document}